% !TeX program = lualatex
% Copyright (C) 2026 IEEE UFG Computer Society contributors.
%
% This work may be distributed and/or modified under the conditions of the
% LaTeX Project Public License, either version 1.3c of this license or (at your
% option) any later version. The latest version is available at
% https://www.latex-project.org/lppl.txt.
%
% This work has the LPPL maintenance status `maintained'. The Current
% Maintainer is the IEEE UFG Computer Society. The work consists of the files
% listed in MANIFEST.md. See LICENSE for details and artwork exclusions.

% GUIA DE EDIÇÃO
% Use uma mensagem principal por quadro e prefira evidências visuais a blocos
% longos de texto. Para apresentações curtas, use
% \usetheme[sectionpages=false]{ieeeufgcs}.
\documentclass[aspectratio=169,11pt]{beamer}
\usetheme{ieeeufgcs}
\usepackage[brazil]{babel}

% Para usar referências bibliográficas, renomeie
% references.bib.example para references.bib e ative as três linhas abaixo.
% \usepackage{csquotes}
% \usepackage[backend=biber,style=numeric]{biblatex}
% \addbibresource{references.bib}

% COMANDOS ÚTEIS
% Imagem: \includegraphics[width=\linewidth]{img/arquivo.png}
% Referência: \cite{chave-da-referencia}
% Destaque numérico: \IEEECSMetric{42}{Descrição breve}

\title[Título breve]{Título da apresentação}
\subtitle{Mensagem principal ou contexto}
\author{Nome da pessoa ou equipe}
\institute{Capítulo Estudantil da IEEE Computer Society na UFG}
\date{Evento ou ocasião · \today}

\begin{document}

\begin{frame}[plain,noframenumbering]
  \titlepage
\end{frame}

\section{Contexto}

\begin{frame}{Contexto e objetivo}
  \begin{columns}[T,onlytextwidth]
    \begin{column}{.48\textwidth}
      {\IEEECSHeadingFont\bfseries\color{IEEEBlue}Contexto\par}
      \medskip
      Apresente as informações necessárias para situar o tema, o público e o
      período abordado.
    \end{column}
    \begin{column}{.48\textwidth}
      {\IEEECSHeadingFont\bfseries\color{IEEEBlue}Objetivo\par}
      \medskip
      Declare o propósito da apresentação e o resultado esperado deste
      encontro.
    \end{column}
  \end{columns}
\end{frame}

\begin{frame}{Evidências}
  \framesubtitle{Indicadores diretamente relacionados ao tema}
  \begin{columns}[T,onlytextwidth]
    \begin{column}{.30\textwidth}
      \IEEECSMetric{42}{Participantes}
    \end{column}
    \begin{column}{.30\textwidth}
      \IEEECSMetric{3×}{Variação no período}
    \end{column}
    \begin{column}{.30\textwidth}
      \IEEECSMetric{87\%}{Índice de conclusão}
    \end{column}
  \end{columns}
  \vfill
  Registre a fonte e o período de referência dos dados apresentados.
\end{frame}

\section{Proposta}

\begin{frame}{Proposta}
  \begin{columns}[c,onlytextwidth]
    \begin{column}{.61\textwidth}
      {\IEEECSHeadingFont\bfseries\Large
        Apresente a proposta em uma frase direta.\par}
      \medskip
      Delimite o escopo, a relação com o contexto e as condições necessárias
      para execução.
    \end{column}
    \begin{column}{.33\textwidth}
      {\IEEECSHeadingFont\bfseries\color{IEEEBlue}Resultado\par}
      \medskip
      Indique a principal entrega ou mudança esperada.
    \end{column}
  \end{columns}
\end{frame}

\begin{frame}{Como funciona}
  \begin{columns}[T,onlytextwidth]
    \begin{column}{.30\textwidth}
      {\IEEECSHeadingFont\bfseries\color{IEEEBlue}01\par}
      \medskip
      \textbf{Preparação}\par
      Confirme responsáveis, recursos e critérios.
    \end{column}
    \begin{column}{.30\textwidth}
      {\IEEECSHeadingFont\bfseries\color{IEEEBlue}02\par}
      \medskip
      \textbf{Execução}\par
      Realize as atividades e registre as evidências.
    \end{column}
    \begin{column}{.30\textwidth}
      {\IEEECSHeadingFont\bfseries\color{IEEEBlue}03\par}
      \medskip
      \textbf{Avaliação}\par
      Compare resultados e defina os ajustes seguintes.
    \end{column}
  \end{columns}
\end{frame}

\begin{frame}{Responsabilidades e recursos}
  \begin{columns}[T,onlytextwidth]
    \begin{column}{.48\textwidth}
      {\IEEECSHeadingFont\bfseries\color{IEEEBlue}Quem participa\par}
      \begin{itemize}
        \item coordenação e tomada de decisão;
        \item execução técnica ou operacional;
        \item apoio institucional e comunicação.
      \end{itemize}
    \end{column}
    \begin{column}{.48\textwidth}
      {\IEEECSHeadingFont\bfseries\color{IEEEBlue}Recursos\par}
      \begin{itemize}
        \item tempo e disponibilidade;
        \item infraestrutura e orçamento;
        \item acessos, dados e autorizações.
      \end{itemize}
    \end{column}
  \end{columns}
\end{frame}

\section{Encaminhamentos}

\begin{frame}{Resultados esperados}
  \begin{itemize}
    \item associe cada resultado a uma evidência observável;
    \item diferencie entregas imediatas de efeitos de longo prazo;
    \item registre limites, riscos e condições de validade;
    \item indique quando os resultados serão avaliados.
  \end{itemize}
  \vfill
  {\IEEECSHeadingFont\bfseries\color{IEEEBlue}Critério de sucesso\par}
  \smallskip
  Escreva uma condição objetiva para verificar o resultado.
\end{frame}

\begin{frame}{Próximos passos}
  \begin{enumerate}
    \item \textbf{Decisão:} diga exatamente o que precisa ser aprovado.
    \item \textbf{Responsável:} identifique quem inicia cada ação.
    \item \textbf{Prazo:} indique a primeira data verificável.
  \end{enumerate}
  \vfill
  {\IEEECSHeadingFont\bfseries\color{IEEEBlue}Encaminhamento\par}
  \smallskip
  Registre a decisão solicitada ao público.
\end{frame}

% Ative este quadro depois de configurar o arquivo references.bib.
% \begin{frame}[allowframebreaks]{Referências}
%   \printbibliography
% \end{frame}

{
  \setbeamercolor{background canvas}{bg=IEEEDarkBlue}
  \begin{frame}[plain]
    \usebeamercolor[fg]{ieeecs closing}
    \centering
    \vfill
    {\IEEECSHeadingFont\bfseries\LARGE
      Retome aqui a única mensagem que o público deve guardar.\par}
    \vspace{6mm}
    Nome ou equipe · contato@example.org
    \vfill
    \IEEECSLogo[variant=orange-white,width=34mm]
  \end{frame}
}

\end{document}
